% STATUS: PHASE 1C AUDIT WITH SPECIFIC ACCESS GAPS; NO PROOFS.
\section{Targeted priority closure}
\label{sec:priority-closure}
\textbf{Decision, 24 September 2026: exact-model priority is not closed.}
The 2021--2024 publication chain is identified, but its complete main texts
have not all been inspected. The publisher's 2024 supplement was obtained and
read in full. This report separates inspected primary content from metadata
and inaccessible content. No canonical equation or proof fragment is changed.
No project theorem is proved and no model computation is performed.

\subsection{Bounded source chain and bibliographic records}
The following records retain the Phase 1B source identifiers. Publication
page ranges are bibliographic locators, not evidence of body inspection.
The publisher's 2024 reference list identifies 2021, 2022 and 2023 as
references 37, 38 and 39, respectively. This is the verified backward link;
the 2024 paper is a verified forward citation of each earlier paper. No
complete forward-citation census is claimed.
\begin{enumerate}[leftmargin=*]
\item \textbf{S13.} Haipeng Lin, Wenxin Shi and Jing Han,
\emph{Analysis of Enzymatic Regulatory Five-Element Model},
2021 40th Chinese Control Conference (CCC), IEEE, pp.~6410--6415.
\href{https://doi.org/10.23919/CCC52363.2021.9550151}{DOI: \nolinkurl{10.23919/CCC52363.2021.9550151}}.
The proceedings contents and 2024 reference 37 confirm the record.
Full main text unavailable in this pass; exact equation and theorem locations
are unresolved, not absent.
\item \textbf{S15.} Haipeng Lin, Wenxin Shi and Jing Han,
\emph{Comparison of Enzymatic Regulatory Network Topologies},
2022 41st Chinese Control Conference (CCC), IEEE, pp.~5735--5740.
\href{https://doi.org/10.23919/CCC55666.2022.9901543}{DOI: \nolinkurl{10.23919/CCC55666.2022.9901543}}.
Proceedings contents and 2024 reference 38 identify it. Main text unavailable.
An author-upload search result under this title contains a 2024 uncorrected
proof headed with the 2024 DOI. It cannot be treated as the 2022 paper.
\item \textbf{S16.} Haipeng Lin and Jing Han,
\emph{Comparison of Multi-Node Networks on Biological Properties},
2023 42nd Chinese Control Conference (CCC), IEEE, pp.~6823--6828.
\href{https://doi.org/10.23919/CCC58697.2023.10239797}{DOI: \nolinkurl{10.23919/CCC58697.2023.10239797}}.
Proceedings contents and 2024 reference 39 identify it. Main text unavailable.
\item \textbf{S14.} Haipeng Lin and Jing Han,
\emph{What is special of ``five'' in biological regulatory networks?},
Nonlinear Dynamics 112(9) (2024), pp.~7477--7498.
\href{https://doi.org/10.1007/s11071-024-09374-5}{DOI: \nolinkurl{10.1007/s11071-024-09374-5}}.
The publisher supplies abstract/references and a subscription preview;
main text unavailable. Its eight-page supplementary PDF is fully inspected.
The download anchor labelled supplementary information points to a 7z archive,
despite the page describing a small PDF. Actual archive and PDF hashes are
recorded in the access log. Author-upload indexed fragments are partial only.
\end{enumerate}

Two narrowly connected leads were identified, without opening another topic.
\textbf{C1-L1:} Haipeng Lin, \emph{Research on Multi-subsystem
Generation/Restriction Coupled Models} (translated title), master's thesis,
University of Chinese Academy of Sciences, Beijing, 2021. The
\href{https://lsc.amss.cas.cn/rcpy/bfbys/202601/t20260122_821421.html}{official laboratory's 2021 graduate list},
entry 18, records the author, title, adviser Jing Han and July 2021 degree;
the thesis body was not obtained. \textbf{C1-L2:} Feng Xiong and Weibo Zhang,
\emph{Analysis of ``Geng Gui Geng Jian'' of the Five Elements in Huangdi
Neijing} (translated title), Acta Chinese Medicine 40(10) (2025),
pp.~2098--2103,
\href{https://doi.org/10.16368/j.issn.1674-8999.2025.10.329}{DOI: \nolinkurl{10.16368/j.issn.1674-8999.2025.10.329}}.
The journal's abstract and references 4, 5 and 7 link S13, C1-L1 and S14.
Its embedded full-text viewer could not be read. These are equation-retrieval
leads, not evidence that their models equal DM.

\subsection{Coverage: what has and has not been inspected}
\textbf{U} means not established from the available body; it never means no.
\textbf{Y} is affirmative inspected evidence; \textbf{N(sup)} means absent
from the inspected supplement only. \textbf{T} means a title-level indication,
not verification of an equation or theorem. Each entry's locator follows in
the notes. ``Cyclic'' describes a graph; ``circulant'' properly describes a
matrix and is not asserted merely from the word symmetric.

\begin{longtable}{>{\raggedright\arraybackslash}p{0.37\textwidth}>{\raggedright\arraybackslash}p{0.105\textwidth}>{\raggedright\arraybackslash}p{0.105\textwidth}>{\raggedright\arraybackslash}p{0.105\textwidth}>{\raggedright\arraybackslash}p{0.19\textwidth}}
\toprule Required field & 2021 & 2022 & 2023 & 2024 main/supplement \\
\midrule\endhead
Exact main equations & U & U & U & U / limits below \\
State count & Five (T) & U & Multi-node (T) & $n\le9$ abstract; S1/S4 \\
Graph topology & Five-element (T) & U & U & Step-one/step-two, S5 \\
Nonlinear interaction law & Enzymatic (T) & Enzymatic (T) & U & Enzymatic abstract; limits S4 \\
Cyclic graph / circulant Jacobian & U/U & U/U & U/U & Cyclic examples S5 / U \\
Five elements, phases or organs & Five-element (T) & U & U & Y, abstract \\
Equilibrium derivations & U & U & U & U / Y, S4 \\
Arbitrary-$n$ theorem & U & U & U & U / N(sup) \\
Fourier eigenvalues & U & U & U & U / N(sup) \\
Exact wave-number selection & U & U & U & U / N(sup) \\
Hopf analysis & U & U & U & S3 cites main Fig.~12; full test U \\
First Lyapunov coefficient & U & U & U & U / N(sup) \\
Nonlinear-law comparison & U & U & U & Limiting-law analysis S4; full inventory U \\
Canonical MB versus DM comparison & U & U & U & U / N(sup) \\
Exact canonical DM present & U & U & U & U / N(sup) \\
\bottomrule
\end{longtable}

\paragraph{Exact inspected supplemental equations and locations.}
S14 supplement S4, p.~4, Eq.~(D2) gives, in compact indexing of its displayed
six equations,
\[
 \dot x_2=ax_7-(b/k)x_1x_2,\qquad
 \dot x_i=ax_{i-2}-(b/k)x_{i-1}x_i,\quad i=3,\ldots,7.
\]
For its nine-node topology, p.~5, Eq.~(D7) gives
\[
 \dot x_2=ax_9-(b/k)x_1x_2,\qquad
 \dot x_i=ax_{i-2}-(b/k)x_{i-1}x_i,\quad i=3,\ldots,9.
\]
These are the source's saturated/linear limiting equations, with the driven
node's equation omitted there; do not reconstruct it. Equation (D5), p.~5,
is an equilibrium system, not an ODE:
\[
 ax_1=(b/k)x_9x_2,\qquad
 ax_{i-1}=(b/k)x_{i-2}x_i,\quad i=3,\ldots,9.
\]
The source sets $A=b/(ak)$ and derives equilibrium relations in (D3)--(D6)
and (D8)--(D9), pp.~4--5. These displayed laws are not either canonical law.
Their existence does not clear the inaccessible main equations.

\paragraph{Overlap locations and limits.}
Supplement S2, pp.~2--3, Definition 1, describes signed connections;
S5, pp.~5--6, Fig.~S5, treats step-one activation/step-two restriction and
the reversed assignment. S3, pp.~3--4, discusses parameter variation near
the main paper's Fig.~12 Hopf region. S4 contains equilibrium/adaptation
calculations and limiting laws. Thus topology, enzymatic motivation and
limiting-law analysis are prior content; the supplement supplies no arbitrary-$n$
selection or Lyapunov theorem. Main-text absence remains untested.

\subsection{Lorenz--96: version-specific mathematical comparison}
All page numbers in this subsection refer to the supplied accepted
manuscript (arXiv version 1704.05442v3; 55 pages), dated 21 December 2017, not to journal pages
38--60. The source is van Kekem and Sterk (2018)\cite{vanKekemSterk2018}.
This is an audit of published results, not a derivation of our proposed results.

\paragraph{Source objects used in the comparison.}
Equation (1), p.~1, is
\[
 \dot x_j=x_{j-1}(x_{j+1}-x_{j-2})-x_j+F.
\]
At $x_F=(F,\ldots,F)$, Eqs.~(3)--(6), pp.~8--9, give
$\lambda_j=-1+Ff(j,n)+iFg(j,n)$, where
\[
 f(j,n)=\cos(2\pi j/n)-\cos(4\pi j/n),\qquad
 g(j,n)=-\sin(2\pi j/n)-\sin(4\pi j/n).
\]
The paper explicitly writes $f=-2z^2+z+1$, $z=\cos(2\pi j/n)$,
in the proof of Theorem 2, p.~37. This is the closest elementary selection
mechanism, not just a thematic resemblance. Proposition 5, pp.~11--12,
selects $\arg\max_{0<l<n/3}f(l,n)$ for the first positive-$F$ instability;
its proof is in \S4.4, pp.~46--47. Lemma 1, p.~10, supplies
$F_H=1/f(l,n)$ and transversality for $n\ge4$, $0<l<n/2$, $l\ne n/3$.
Theorem 2, p.~10, Eq.~(7), characterizes distinct-pair coincident crossings by
\[
 \cos(2\pi l_1/n)+\cos(2\pi l_2/n)=\tfrac12.
\]
This is a spectral coincidence criterion; a nondegenerate nonlinear
codimension-two unfolding requires additional conditions.

Proposition 6, p.~13, states
\[
 \lim_{n\to\infty}\frac{n}{l_1(n)}=
 \frac{2\pi}{\arccos(1/4)},\qquad
 T_\infty=2\pi\tan\!\left(\tfrac12\arccos(1/4)\right).
\]
Theorem 4, p.~11, and Eq.~(36), p.~44, give an explicit $\ell_1(l,n)$
and classify its sign. The numerator controlling that sign is
\[
 L(y)=5\cos y+8\cos(2y)-2\cos(3y)-8,
 \quad y=2\pi l/n
\]
(Eq.~(37), p.~45). The unique zero $y_0\in(0,\pi)$ separates positive
and negative sign on the stated mode range; \S4.3.2, pp.~44--46,
establishes the sign, and Proposition 5 gives supercriticality of a simple
first positive-$F$ Hopf. Our DM sign is not inherited from this result.

\paragraph{Two-parameter overlap must not be omitted.}
Equation (9), p.~14, adds $G(x_{j-1}-2x_j+x_{j+1})$ to the source model.
Equation (10) and Lemma 8/Eq.~(11), p.~15, give its circulant spectrum
and linear crossing curves in $(F,G)$. Section 2.4.2, pp.~16--18, treats
the $n=12$ unfolding. Thus the source also treats competing crossing curves;
it is inaccurate to dismiss its switching relevance because $F$ alone
rescales the one-parameter growth function. It does not compare our
rational DM law with our bilinear comparator.

\clearpage
\begin{landscape}
\subsubsection*{Theorem-by-theorem comparison (eight required objects)}
\small
\setlength{\tabcolsep}{3pt}
\renewcommand{\arraystretch}{1.15}
\begin{longtable}{>{\raggedright\arraybackslash}p{0.105\linewidth}>{\raggedright\arraybackslash}p{0.16\linewidth}>{\raggedright\arraybackslash}p{0.115\linewidth}>{\raggedright\arraybackslash}p{0.13\linewidth}>{\raggedright\arraybackslash}p{0.155\linewidth}>{\raggedright\arraybackslash}p{0.10\linewidth}>{\raggedright\arraybackslash}p{0.155\linewidth}}
\toprule Mathematical object & Van Kekem--Sterk result & Assumptions & Proof method & Our proposed result & Relation & Extra content needed \\
\midrule\endhead
Finite-$n$ selection & Prop.~5, pp.~11--12; quadratic explicitly in \S4.2 p.~37; proof \S4.4 pp.~46--47. & $n\ge4$, first crossing for $F>0$; eligible $0<l<n/3$; pair ties possible. & Circulant spectrum, discrete maximization of $\cos\theta-\cos2\theta$, interval comparison. & P05: exact $q=2$ selection throughout a stated positive-parameter/branch domain. & Analogous; bare rule routine. & Realizable cells, ties, endpoints and first-instability meaning; use in P08 contrast. \\
\addlinespace
Large-$n$ wave number & Prop.~6, p.~13: $n/l_1\to2\pi/\arccos(1/4)$; period limit. & Same first positive-$F$ selection; fixed symbol as $n$ increases. & Dense angular sampling of the continuous maximum. & P06: limit for a fixed DM family, if informative. & Analogous; generic argument routine. & A meaningful comparison or controlled finite-size effect; define the branch/family and degeneracies. \\
\addlinespace
Critical-mode switching & Thm.~2/Eq.~(7), p.~10; proof pp.~37--38. Extended system Eq.~(9), p.~14; crossing lines (11), p.~15; $n=12$ case pp.~16--18. & Distinct eligible pairs; $n\ge4$; extra nonlinear conditions for a Hopf--Hopf unfolding. & Equal values of a quadratic; equality/intersections of crossing curves in $(F,G)$. & P05/P08: parameter boundaries between first modes in DM and MB. & Analogous. & Show boundaries are admissible and reached on the stated path; certify first crossing and all other modes; classify actual law contrast. \\
\addlinespace
Hopf conditions & Lemma~1 p.~10, Thm.~4 p.~11; proofs \S4.1 and \S4.3, pp.~33--46. & $n\ge4$, $0<l<n/2$, $l\ne n/3$, unique imaginary pair, transversality and nonzero $\ell_1$. & Explicit spectrum, crossing derivative, standard Hopf theorem and coefficient calculation. & P11--P13: first symmetry-breaking Hopf on a regular positive DM branch. & Analogous; theorem application known. & Nonempty certified parameter regime, moving-equilibrium derivative, nonzero frequency, all other modes stable and $\ell_1\ne0$. \\
\addlinespace
First Lyapunov coefficient & Thm.~4 p.~11; Eqs.~(23)--(36), pp.~39--44; zero- and doubled-mode contractions in pp.~42--43. & Simple Hopf; specified unitary normalization; quadratic vector field, cubic tensor zero. & Normality gives paired eigenvectors; Fourier products remove matrix inverses; scalar resolvents. & P13: exact coefficient for rational DM, including cubic derivatives. & Analogous; exact expression model-specific. & Correct derivatives, tensor factors, harmonic coincidences and resolvent domains. Nonzero cubic tensor alone is not methodological novelty. \\
\addlinespace
Sign/criticality & Thm.~4 p.~11; \S4.3.2 pp.~44--46; Prop.~5 pp.~11--12 and proof p.~47. & Eligible simple Hopf; sign intervals relative to $y_0$; simple first positive-$F$ crossing is supercritical. & Positive-factor separation; trigonometric-polynomial derivative and unique zero. & P13: nontrivial analytic DM sign region with a selection or DM--MB consequence. & Analogous goal; proposed sign law unresolved. & Certify a nonempty region and full claimed size/mode range; exclude zeros or analyze only to justified order. \\
\addlinespace
Rotating or traveling wave & \S2.3 pp.~12--13, Eq.~(8): leading periodic wave, wavelength, period and direction. & Local Hopf branch; Fourier mode; leading-order amplitude approximation. & Interpret the critical Fourier eigenvector and oscillatory phase. & P14: exact spatiotemporal symmetry and isotropy of the nonlinear branch. & Analogous; exact symmetry also uses known equivariant theory. & Distinguish exact branch symmetry from linear approximation; record phase sign, stabilizer and branch stability conditions. \\
\addlinespace
Two coupling laws & Eq.~(1) p.~1 versus diffusive extension Eq.~(9) p.~14; no DM--MB comparison in this inspected source. & Original quadratic advection retained; new linear diffusion coefficient $G$. & Compare spectra/crossing curves of original and extended systems. & P08: same-coefficient DM versus MB with independently specified positive branches. & Genuinely different comparison object; novelty conditional. & A rigorous qualitative contrast on a nonempty common admissible set. Source-chain clearance and mathematical verification remain separate. \\
\bottomrule
\end{longtable}
\end{landscape}

\subsection{Decision on the q=2 theorem}
For the \emph{bare quadratic maximizer}, the correct choice is
\textbf{routine adaptation}. It is not yet a genuinely stronger theorem.
For the strengthened all-admissible-parameter selection, switching and
matched-comparison target, the correct planning description is
\textbf{structurally analogous but model-specific}. Changing coefficients
and substituting an equilibrium is not enough: realizability and qualitative
consequences are the intended added content. None is proved here.

\subsection{Final claim verdicts and frozen paper roles}
The following labels concern novelty only. Each rating remains subject to the
stated source-access limitation; no rating establishes correctness or priority.
\begin{longtable}{>{\raggedright\arraybackslash}p{0.32\textwidth}>{\raggedright\arraybackslash}p{0.32\textwidth}>{\raggedright\arraybackslash}p{0.25\textwidth}}
\toprule Claim & Verdict & Proof owner \\ \midrule\endhead
Arbitrary-(n,q) spectrum & KNOWN / ROUTINE & P04 \\
$q=2$ exact finite-n selection & SUPPORTING SPECIALIZATION & P05 \\
Mode-switch boundaries & POSSIBLE NOVEL CONTRIBUTION & P05, P08 \\
Large-n selected-wave-number limit & KNOWN / ROUTINE & P06 \\
Matched DM--MB comparison & STRONG NOVELTY CANDIDATE & P08; P11--P13 if criticality is compared \\
Rigorous Hopf theorem & SUPPORTING SPECIALIZATION & P11--P13 \\
First Lyapunov coefficient reduction & SUPPORTING SPECIALIZATION & P13 \\
Sign/criticality result for $\ell_1$ & UNRESOLVED & P13 \\
Rotating-wave conclusion & SUPPORTING SPECIALIZATION & P14 \\
\bottomrule\end{longtable}
The comparison rating concerns a substantive contrast, not two formulas.
The sign verdict remains unresolved because no valid region has been stated
and established. The bare large-$n$ limit and general spectrum remain routine.
For detailed conditions attached to each verdict see
\texttt{final\_novelty\_verdict.md}.

The supporting-theory locators retained from Phase 1B are Gray
\cite{Gray2006}, \S3.1, printed pp.~32--34, Eqs.~(3.7)--(3.13),
and \S3.2, Theorem 3.1, printed pp.~34--35 (circulant structure);
Pais et al.\cite{PaisCaicedoLeonard2012}, supplied 30-page version,
Lemma 4.6, p.~13, Appendix D, pp.~26--28, Eqs.~(D.1)--(D.4)
(Fourier Lyapunov reduction including cubic terms); and
Filipski--Golubitsky\cite{FilipskiGolubitsky2010}, \S3, Theorem 3.1,
p.~285, and \S5, Theorem 5.1, p.~288 (abelian Hopf symmetry).
These are precise locations of known machinery, not an assertion that a
realization theorem proves existence in this particular DM vector field.

The frozen package has four intended main results: selection with realizable
switch cells (P05), a matched-law contrast (P08), an assembled first Hopf
classification (P11--P13), and a substantive criticality result (P13).
The coefficient reduction is a proposition/lemma; rotating waves and the
five-node specialization are corollaries. General spectrum, equilibria,
transcritical and fold results support them. Retain P06 only as an informative
large-$n$ corollary. P15--P18 are postponed, including their optional figures.
Detailed standard estimates and long coefficient algebra belong in a supplement.

The primary comparison uses the same positive $(a,k,c,g,h)$, $n,q$ and time
variable, each model evaluated at its own explicitly chosen positive synchronous
branch on the common existence domain. The branch/path and a valid sign region
are future statement gates; this audit does not invent them.

\subsection{Stop decision and release gate}
The requested priority-closure condition is unmet: identification and partial
inspection of the source chain are not full primary-source inspection.
Exact equations unavailable for S13, S15, S16 and the S14 main article must
remain explicitly unresolved. The inspected supplement rules out neither an
exact antecedent elsewhere in its parent article nor a theorem in an earlier
conference version. No universal absence claim follows from this audit.

Stop literature work after this targeted pass. Do not commission a new general
survey. If these named bodies become available, the remaining operation is
their exact equation/theorem comparison against this inventory. Their absence
blocks priority assertions, not later independent mathematical work. All proof
files retain TODO status and the immutable master remains unchanged.
